\section{A focused processor-literature assessment}\label{sec:census}
\subsection{Search scope and the distinction being tested}
The post-Edition-2 discussion suggested that no nontrivial processor-class witness existed. That universal negative is not supported. This assessment follows the cited physical-network work and checks primary reports of non-electronic logic, finite automata, programmed photonics and large-model proposals. It is a focused source assessment as of 10 September 2026, not an exhaustive or preregistered systematic review. The source register states what was read and distinguishes publisher or author abstracts from inspected full-text passages.

The question is not whether a paper uses our notation. It is which part of a proposed realisation obligation its demonstrated apparatus can support. A role-level precedent, a calibrated transformation, a stateful machine, an equivalent pair and a live transfer of a current state are different accomplishments. We supply no new physical experiment in this edition.

\subsection{Physical state and transformations already demonstrated}
Prakash and Gershenfeld's \emph{Microfluidic Bubble Logic} reports engineered fluidic gates, bistable toggle memory and counters \citep{bubble}. The toggle stores a bubble configuration that changes under a subsequent event; the full-text description and the relevant figure were inspected. This is a processor-and-memory precedent, not merely a communication carrier. It is not evidence that ambient turbulence is programmable or that a digital execution was transferred into the apparatus.

Benenson and colleagues report a programmable biomolecular finite automaton whose molecular input, transition rules and final-state output are realised in a specified DNA/enzyme system \citep{benenson}. The primary abstract supports a finite-state computation claim, not an arbitrary general-purpose computer or a digital-to-molecular handover. Both examples undermine the unrestricted assertion that a non-electronic medium can only carry messages. Neither completes the larger host-independence claim.

Wright and colleagues demonstrate trained optical, mechanical and electrical physical networks using hybrid physical--digital physics-aware training \citep{wright}. The primary abstract and publication metadata were checked; remembered parameter counts and precision numbers are not used as measurements. Task learning in these systems is relevant evidence for physical processing, but it is not a demonstration of equivalence to a particular large stateful transformer.

Shen and colleagues' inspected preprint reports training on a conventional computer and programming the resulting transformations into a silicon-photonic system \citep{shen}. The linear computation uses optical interference, while the reported proof of concept implements the nonlinear transformation electronically. Thus ``every physical network must learn all parameters from scratch on its own medium'' is false, but ``a full model ran without electronic support'' would also misstate this example. The relevant experimental and methods passages, including their figures, were inspected. No numerical result from that preprint is silently substituted for a different final-publication version.

\begin{center}\small
\begin{tabular}{@{}>{\raggedright\arraybackslash}p{0.25\textwidth}>{\raggedright\arraybackslash}p{0.31\textwidth}>{\raggedright\arraybackslash}p{0.34\textwidth}@{}}\toprule
Source & Supported role or procedure & Not supplied by the cited observation\\\midrule
Bubble logic \citep{bubble} & Physical stateful logic and memory & Cross-device stateful handover\\
Biomolecular automaton \citep{benenson} & Molecular finite-state computation & Arbitrary transformer implementation\\
Physical neural networks \citep{wright} & Trained physical transformations & Uniform correspondence to a fixed target machine\\
Programmed photonics \citep{shen} & Digital training, optical linear execution, electronic support & Independence from all electronic infrastructure\\
Physical foundation models \citep{pfm} & Proposal and scaling discussion & A demonstrated large-scale realised model\\\bottomrule
\end{tabular}
\end{center}
These verdicts are scoped to the claims inspected. ``Not supplied'' is not a proof that no related experiment exists elsewhere. No row is promoted to a paired migration merely because it can implement a task also executable in software.

\subsection{The large-model question is active, but a proposal is not a device}
A 2026 preprint by Wright, Wang, Onodera and McMahon proposes fixed physical implementations of foundation models and presents back-of-the-envelope scaling calculations, including a nanostructured-glass optical example \citep{pfm}. Its abstract frames major unresolved engineering challenges. That is a relevant research proposal, not a realised trillion-parameter device or a bound certifying a model's full inference trace. Its existence corrects an overly broad literature-negative claim without establishing the ambitious application.

An exact or approximate contract for a large trained model remains substantially stronger than a small finite-state demonstration. This assessment neither certifies such a model on an unconventional processor nor proves a universal scale barrier. The benchmark in Section~\ref{sec:witness} instead isolates statefulness and transfer under a small contract so that these can be tested without a large model's many simultaneous burdens.

\subsection{What the literature finding changes}
The defensible conclusion has three parts. Alternative physical computation is an established experimental subject. The paired preservation of a chosen stateful contract across specified realisations needs its own evidence. The speculative organisation described by the broader motivating scenario is not demonstrated by the papers above. Each statement has a different bearer and evidential burden; none should be used as a rhetorical substitute for the others.
